\section*{Impact Statement}
This work aims to make explanations of language-model internals more reliable, by showing that the cache channel that appears to carry a value or an intervention depends on the prompt under which it is measured. More reliable component-level explanations could help researchers audit and debug models, and could reduce the risk that a mechanism found under one evaluation format is trusted to explain behaviour under another. Our results also caution against over-reading interpretability findings: an attribution validated under one prompt format, or inferred from the model's own read, should not on its own certify behaviour under others. Methods for locating and editing what a model reads from its context could also be used to steer models in unwanted ways; our experiments use public models, templated stories, public reading-comprehension passages, edits built from public material, a released remap and a released state-token patch.
